% Shared geometry for repeated momentum problems. Captions remain at call sites.
\newcommand{\momentumRain}{%
\begin{tikzpicture}[scale=.78]
  \fill[Blue!8] (-1.2,-.65) rectangle (1.2,-.12);
  \draw[solbody] (-1.2,.7)--(-1.2,-.65)--(1.2,-.65)--(1.2,.7);
  \foreach \x in {-.9,-.45,0,.45,.9}{\draw[solflow] (\x,1.55)--(\x,.86);}
  \node[sollabel] at (0,2.05) {$q=\rho S\Delta h/\tau$，雨水水平速度为零};
  \node[sollabel] at (0,-1.15) {$m(t)=M+qt$};
  \draw[solforce] (1.2,-.38)--(2.55,-.38) node[midway,below=3pt,sollabel] {$F$};
  \draw[solflow] (1.6,.52)--(2.8,.52) node[midway,above=3pt,sollabel] {$v$};
\end{tikzpicture}}

\newcommand{\momentumPickup}{%
\begin{tikzpicture}[scale=.8]
  \draw[solbody] (-2,0)--(2.5,0);
  \draw[DeepBlue,line width=1.2pt] (-1.8,0)--(-.15,0)
    .. controls (-.05,0) and (0,.05) .. (0,.15)--(0,2.4);
  \fill[Ink] (0,2.4) circle (2pt);
  \draw[solforce] (0,2.4)--(0,3.3) node[left=3pt,sollabel] {$F$};
  \draw[solforce] (-.65,1.65)--(-.65,.7) node[midway,left=3pt,sollabel] {$\lambda xg$};
  \draw[solflow] (.65,1.1)--(.65,2.1) node[midway,right=3pt,sollabel] {$v$};
  \draw[solgeom] (0,2.4)--(1.8,2.4);
  \draw[<->,draw=TextGray] (1.8,0)--(1.8,2.4) node[midway,right=3pt,sollabel] {$x$};
  \node[sollabel] at (-1.0,-.38) {静止绳段};
  \node[sollabel] at (0,-.95) {$\dot m_{\rm mov}=\lambda v$，新增动量率 $\lambda v^2$};
\end{tikzpicture}}

\newcommand{\momentumFallingRope}{%
\begin{tikzpicture}[scale=.78]
  \draw[solbody] (-.45,3)--(.45,3);
  \foreach \x in {-.4,-.2,0,.2,.4}{\draw[TextGray] (\x,3)--++(.12,.14);}
  \draw[solgeom] (0,3)--(0,2.1);
  \draw[DeepBlue,line width=1.1pt] (0,2.1)--(0,.28)
    .. controls (0,.08) and (-.15,.08) .. (-.3,.08);
  \draw[DeepBlue,line width=1.1pt,domain=-.3:1.6,samples=120]
    plot ({\x},{.08+.08*sin((\x+.3)*720)});
  \draw[solbody] (-.75,0)--(1.9,0);
  \draw[solgeom] (-.7,3)--(0,3) (-.7,2.1)--(0,2.1);
  \draw[<->,draw=TextGray] (-.7,2.1)--(-.7,3) node[midway,left=3pt,sollabel] {$x$};
  \draw[solflow] (.55,1.8)--(.55,.9) node[midway,right=3pt,sollabel] {$v$};
  \node[sollabel,anchor=west] at (.65,2.8) {虚线：上端已下降的距离};
  \node[sollabel,anchor=west] at (1.15,1.5) {$l-x$：仍在下落的绳长};
  \node[sollabel] at (.4,-.4) {秤盘；堆积绳段长度为 $x$};
\end{tikzpicture}}

\newcommand{\momentumRopeLoad}{%
\begin{tikzpicture}[scale=.78]
  \draw[DeepBlue,line width=1.1pt] (0,2)--(0,.28)
    .. controls (0,.08) and (-.15,.08) .. (-.3,.08);
  \draw[DeepBlue,line width=1.1pt,domain=-.3:1.8,samples=120]
    plot ({\x},{.08+.08*sin((\x+.3)*720)});
  \draw[solbody] (-.65,0)--(2.05,0);
  \draw[solflow] (.55,1.7)--(.55,.85) node[midway,right=3pt,sollabel] {$v$};
  \draw[solforce] (.45,0)--(.45,-1.0) node[below=3pt,sollabel] {$\lambda xg$};
  \draw[solforce] (1.6,0)--(1.6,-1.0) node[below=3pt,sollabel] {$\lambda v^2$};
  \node[sollabel,anchor=west,align=left] at (2.7,1.15)
    {秤盘承受的向下作用\\$N=\lambda xg+\lambda v^2$\\第二项来自使绳元停止};
\end{tikzpicture}}

\newcommand{\momentumAccretion}{%
\begin{tikzpicture}[scale=.8]
  \draw[solbody,fill=SoftGray] (-.8,0) rectangle (.8,.9);
  \draw[solforce] (-.2,.45)--(-.2,-.85) node[below=3pt,sollabel] {$mg$};
  \draw[solflow] (1.25,.7)--(1.25,-.25) node[midway,right=3pt,sollabel] {$v$};
  \foreach \x in {-1.3,-.85,-.4,.05,.5,.95,1.4}{\fill[TextGray] (\x,-1.55) circle (1.5pt);}
  \node[sollabel] at (0,-1.95) {地面系中静止的介质};
  \draw[<->,TextGray] (-.8,1.3)--(.8,1.3) node[midway,above=3pt,sollabel] {截面积 $S$};
  \node[sollabel,anchor=west,align=left] at (2.3,.8)
    {$\dot m=\rho Sv$\\$\dd(mv)/\dd t=mg$\\$m\dot v=mg-\rho Sv^2$};
\end{tikzpicture}}

\newcommand{\momentumDroplet}{%
\begin{tikzpicture}[scale=.8]
  \draw[solbody,fill=Blue!8] (0,.5) circle (.7);
  \fill[Ink] (0,.5) circle (1.5pt);
  \draw[solgeom] (0,.5)--(.7,.5) node[midway,above=4pt,sollabel] {$r$};
  \draw[solforce] (0,.5)--(0,-1.0) node[below=3pt,sollabel] {$mg$};
  \draw[solflow] (1.25,.75)--(1.25,-.25) node[midway,right=3pt,sollabel] {$v$};
  \foreach \x in {-1.3,-.85,-.4,.05,.5,.95,1.4}{\fill[TextGray] (\x,-1.65) circle (1.5pt);}
  \node[sollabel] at (0,-2.05) {地面系中静止的雾滴};
  \node[sollabel,anchor=west,align=left] at (2.3,.8)
    {$S=\pi r^2$\\$\dot m=\rho_f\pi r^2v$\\$\dot r=\rho_f v/(4\rho_w)$};
\end{tikzpicture}}
